\usepackage{amsfonts}
\usepackage{amstext}
\usepackage{amsopn}
\usepackage{amsthm}
\usepackage{amssymb}
\usepackage{amsmath}
\usepackage{hyperref}

\begin{document}

%Geometric Resolution of the Residual Flavor Orientation in the Noncommutative Standard Model
\title{On the Limits of Spectral Geometry in Determining Flavor Mixing}

\maketitle

\begin{abstract}
The finite spectral triple formulation of the Standard Model geometrizes fermion mass spectra but does not, in general, fix the relative orientation of generation subspaces. In the lepton sector, we show that whenever the effective neutrino Hermitian operator
\[
H_\nu = m_\nu m_\nu^\dagger
\]
admits a two-dimensional eigenspace, an associated residual $\mathrm{SO}(2)$ freedom arises in the $\mu$-$\tau$ plane. This internal orientation symmetry is invisible to the spectral action and compatible with all axioms of noncommutative geometry, yet it acts nontrivially on the PMNS matrix by leaving the atmospheric mixing angle $\theta_{23}$ undetermined. This mechanism explains the persistence of the $\theta_{23}$ ambiguity across operator-first constructions, seesaw deformations, and basis-independent spectral selection principles.

We prove that, subject to the order-zero, first-order, and reality conditions, the unique minimal relational operator capable of resolving this residual symmetry is a linear combination of sector Hermitians,
\[
A(\alpha,\beta) = \alpha H_e + \beta H_\nu,
\]
defined up to real scaling. This operator selects a preferred relative orientation between the charged-lepton and neutrino sectors without modifying the spectral triple. To quantify this mechanism, we introduce an alignment energy functional that measures the failure of $H_e$ and $H_\nu$ to be simultaneously diagonal within the unfixed subspace. We show that this functional is real-analytic and possesses a unique global minimum, at which the PMNS matrix takes the aligned form
\[
U_{\mathrm{PMNS}} = U_e^\dagger U_\nu R(\theta),
\]
thereby fixing $\theta_{23}$.

The alignment mechanism leaves all fermion masses, mass hierarchies, and CKM parameters unchanged, modifying only the geometric degree of freedom left indeterminate by the spectral triple. These results delineate the boundary of what can be fixed by spectral geometry alone and provide a minimal, fully geometric resolution of the final unfixed lepton mixing parameter in the noncommutative Standard Model.
\end{abstract}


\section{Background}

\begin{quote}
\("\)Within this framework, both gravity and gauge interactions emerge from spectral properties, and the Higgs field is realized as inner fluctuations of $D$.\("\)
\emph{- Alain Connes, Noncommutative Geometry (1994).}
\end{quote}

In this section we summarize the elements of the finite spectral triple relevant for the flavor-alignment analysis developed in later sections. Our aim is not to reproduce the full construction of the noncommutative Standard Model, but to isolate the algebraic and geometric structures that determine fermion masses and mixings.

\subsection{The Finite Spectral Triple of the Standard Model}

The internal (finite) sector of the Standard Model is encoded by the spectral triple
\[
(\mathcal A_F,\mathcal H_F,D_F,J_F,\gamma_F),
\]
where:
\begin{itemize}
\item $\mathcal A_F=\mathbb C \oplus \mathbb H \oplus M_3(\mathbb C)$ acts on internal gauge and flavor indices;
\item $\mathcal H_F$ is the finite-dimensional Hilbert space of quark and lepton multiplets, including chiralities and particle-antiparticle sectors;
\item $D_F$ encodes the Yukawa couplings and Majorana mass terms in block-diagonal form;
\item $J_F$ is the real structure implementing charge conjugation;
\item $\gamma_F$ is the internal grading corresponding to chirality.
\end{itemize}

The order-zero condition
\[
[\pi(a),J_F \pi(b)^* J_F^{-1}] = 0,\qquad \forall a,b\in\mathcal A_F,
\]
forces $\mathcal A_F$ to act diagonally on generation space, so the fermionic generations appear as multiplicities in the representation. In particular, the algebra does not mix generations, a fact crucial for the analysis that follows.

The first-order condition
\[
[[D_F,\pi(a)],J_F\pi(b)^*J_F^{-1}] = 0
\]
restricts the allowed form of the Dirac operator and ensures that commutators between geometric one-forms respect a Leibniz rule.

Throughout the paper, we work entirely within this fixed finite spectral triple; no modification of $\mathcal A_F$, $\mathcal H_F$, or $D_F$ will be introduced.

``A spectral triple $(\mathcal A,\mathcal H,D)$ consists of an involutive algebra $\mathcal A$ represented on a Hilbert space $\mathcal H$ together with a self-adjoint operator $D$ with compact resolvent such that $[D,a]$ is bounded for all $a\in\mathcal A$. This data generalizes the notion of a Riemannian spin manifold: $\mathcal A$ plays the role of functions on the manifold, $\mathcal H$ the space of spinors, and $D$ the Dirac operator.'' - Alain Connes, \emph{Noncommutative Geometry} (1994).

\subsection{Sector Hermitian Operators and Mixing Matrices}

In each fermion sector $(f\in\{e,\nu,u,d\})$, the Dirac operator contains a Yukawa block
\[
Y_f : \mathbb C^3 \longrightarrow \mathbb C^3,
\]
which couples left- and right-handed fields. Physical masses and mixing angles depend not on the Yukawa matrices themselves, but on the Hermitian combinations
\[
H_f = Y_f Y_f^\dagger.
\]

The eigenvalue decompositions
\[
H_f = U_f \Lambda_f U_f^\dagger
\]
encode:
\begin{itemize}
\item the fermion mass spectra via $\Lambda_f = \mathrm{diag}(m_{f_1}^2,m_{f_2}^2,m_{f_3}^2)$;
\item the fermion mixing matrices via relative orientations:
\[
U_{\mathrm{CKM}} = U_u^\dagger U_d,\qquad
U_{\mathrm{PMNS}} = U_e^\dagger U_\nu.
\]
\end{itemize}

It is therefore the relative geometry of the Hermitian operators $H_e$ and $H_\nu$-not their individual spectra-that determines neutrino mixing.

\subsection{Degeneracy in the Neutrino Sector}

A key structural feature of the noncommutative Standard Model, and the starting point of our analysis, is the presence of a two-dimensional degeneracy in the neutrino Hermitian operator $H_\nu$. In many formulations of the internal Dirac operator, or in seesaw scenarios with rank-one or nearly rank-one heavy Majorana masses, one finds
\[
H_\nu v = \lambda v \qquad \text{for all } v\in E,
\]
where $E\subset\mathbb C^3$ is a two-dimensional subspace. It suffices that $H_\nu$ has an approximate or exact two-dimensional eigenspace stable under admissible perturbations.

This degeneracy implies that the restriction of $H_\nu$ to $E$ is proportional to the identity:
\[
H_\nu|_E = \lambda \mathbf 1_E.
\]

Consequently, any rotation
\[
R(\theta)\in \mathrm{SO}(2)
\]
acting within $E$ leaves $H_\nu$ invariant:
\[
R(\theta)^\dagger H_\nu R(\theta) = H_\nu.
\]

Such rotations do not violate the order-zero or first-order conditions, nor do they affect the spectral action, which depends only on the eigenvalues of $H_\nu$. Thus the finite spectral triple itself cannot distinguish between different choices of basis in $E$.

This residual $\mathrm{SO}(2)$ symmetry is the geometric origin of the indeterminacy of the atmospheric mixing angle $\theta_{23}$, as shown in Theorem~1.

\subsection{Geometric Interpretation of Mixing Angles}

Mixing angles arise from the relative orientation of the charged-lepton and neutrino eigenbases. If
\[
H_e = U_e \Lambda_e U_e^\dagger,\qquad
H_\nu = U_\nu \Lambda_\nu U_\nu^\dagger,
\]
then
\[
U_{\mathrm{PMNS}} = U_e^\dagger U_\nu.
\]

A rotation $R(\theta)$ in the degenerate neutrino subspace induces
\[
U_\nu \longmapsto U_\nu R(\theta),
\]
and therefore
\[
U_{\mathrm{PMNS}} \longmapsto U_{\mathrm{PMNS}} R(\theta).
\]

Since $R(\theta)$ acts only in the $\mu$-$\tau$ plane, its effect is confined to the atmospheric mixing angle $\theta_{23}$. The solar and reactor angles $(\theta_{12},\theta_{13})$ remain unchanged.

This geometric structure isolates a single physical degree of freedom-the $\mu$-$\tau$ rotation-whose value must be fixed in order to determine $\theta_{23}$.

The remainder of the paper addresses:
\begin{itemize}
\item how this degree of freedom arises (Section~3),
\item how it may be fixed by a minimal, NCG-admissible operator (Section~4), and
\item how this mechanism yields a unique atmospheric angle (Section~6).
\end{itemize}

It is important to emphasize that spectral triples encode geometric information through spectral invariants. While this suffices to determine eigenvalue data, the axioms do not, by themselves, determine a preferred eigenbasis in the presence of degeneracies. Consequently, any physical quantity depending on the relative orientation of eigenspaces-such as fermion mixing matrices-may remain undetermined by the spectral action, even when all axioms are satisfied.

\begin{quote}
\("\)In noncommutative geometry, the geometry of a space is completely encoded in the spectral properties of the operator $D$ acting in $\mathcal H$. Two spaces are isometric if and only if their associated spectral triples are unitarily equivalent.\("\)

\emph{- Paraphrase of the core philosophy of NCG \`a la Connes (1994).}
\end{quote}


\section{The Residual Orientation Obstruction}

In this section we identify and characterize the geometric origin of the undetermined atmospheric mixing angle. The key observation is that whenever the effective neutrino Hermitian operator admits a two-dimensional eigenspace within generation space, an associated internal $\mathrm{SO}(2)$ symmetry arises. This symmetry is left unfixed by the axioms of the finite spectral triple and is not detected by the spectral action. We show that it corresponds precisely to the freedom to rotate the $\mu$-$\tau$ plane, and that it acts nontrivially on the PMNS matrix while leaving all geometric data of the spectral triple unchanged.

We emphasize that we do not claim that the axioms of noncommutative geometry enforce an exact two-fold degeneracy in the neutrino sector. Rather, the observation is structural: if the effective neutrino Hermitian operator exhibits a two-dimensional eigenspace-whether exact or approximate-then the associated rotational freedom leaves all spectral and geometric data invariant. In this case, the atmospheric mixing angle corresponds precisely to this undetermined internal orientation. The obstruction identified here is therefore not tied to a specific model, parameter choice, or dynamical assumption, but to the presence of a neutrino multiplicity within the internal geometry.

\subsection{Degeneracy of $H_\nu$ and its Consequences}

Let
\[
H_\nu = m_\nu m_\nu^\dagger
\]
be the Hermitian neutrino sector operator. As discussed in Section~2.3, the structure of the internal Dirac operator implies the existence of a two-dimensional subspace
\[
E \subset \mathbb C^3
\]
on which $H_\nu$ is proportional to the identity:
\begin{equation}
H_\nu|_E = \lambda \mathbf 1_E,
\qquad \dim E = 2.
\end{equation}

This degeneracy means that for any rotation
\[
R(\theta) \in \mathrm{SO}(2),
\]
acting unitarily within $E$, one has
\[
R(\theta)^\dagger H_\nu R(\theta) = H_\nu.
\]

Since the representation $\pi(a)$ of $\mathcal A_F$ is diagonal in generation space, the order-zero condition
\[
[\pi(a),J_F\pi(b)^*J_F^{-1}] = 0
\]
is unaffected by such rotations. Likewise, the first-order condition remains untouched. Moreover, the spectral action depends only on the spectrum of $D_F$, which is invariant under conjugation by $R(\theta)$.

Thus every choice of rotation $R(\theta)$ yields a valid internal Dirac operator with identical geometric and spectral data. The freedom to choose the orientation of $E$ therefore remains as an undetermined degree of freedom in the spectral triple.

``The incorporation of right-handed neutrinos and a see-saw mechanism arises naturally within the modified spectral triple, solving previous fermion-doubling issues.'' - Summary from Connes's work on neutrino mixing in NCG.

\subsection{Action on the PMNS Matrix}

Let
\[
H_e = Y_e Y_e^\dagger
\]
be the charged-lepton Hermitian operator, and write the diagonalizations
\[
H_e = U_e \Lambda_e U_e^\dagger,\qquad
H_\nu = U_\nu \Lambda_\nu U_\nu^\dagger.
\]

The PMNS matrix is constructed from their relative orientation:
\[
U_{\mathrm{PMNS}} = U_e^\dagger U_\nu.
\]

A rotation $R(\theta)$ in the degenerate neutrino subspace induces the transformation
\[
U_\nu \longmapsto U_\nu R(\theta)
\quad\Longrightarrow\quad
U_{\mathrm{PMNS}} \longmapsto U_{\mathrm{PMNS}} R(\theta).
\]

Since $R(\theta)$ acts trivially on the electron row and mixes only the $\mu$ and $\tau$ indices, its effect on the mixing angles is:
\begin{itemize}
\item $\theta_{12}$ unchanged,
\item $\theta_{13}$ unchanged,
\item $\theta_{23}$ shifts by $\theta$.
\end{itemize}

Thus the residual $\mathrm{SO}(2)$ symmetry corresponds exactly to the freedom to rotate the atmospheric mixing angle. This is the structural reason the spectral triple cannot fix $\theta_{23}$.

``In the noncommutative geometry formulation of the standard model, neutrino mixing is a structural outcome of the doubling of the algebra and the spectral properties of the finite geometry. This indicates that the mismatch between flavor and mass bases is tied to underlying algebraic structure rather than arbitrary parameter choices.'' - Devastato, Lizzi \& Martinetti 2014.

\subsection{The Residual Orientation Obstruction}

We now formalize this observation.

\begin{theorem}[Residual Orientation Obstruction]
Let $E\subset\mathbb C^3$ be the two-dimensional degenerate eigenspace of $H_\nu$. For every rotation $R(\theta)\in\mathrm{SO}(2)$ acting in $E$:

\begin{enumerate}
\item \emph{Spectral triple invariance.}
\[
R(\theta)^\dagger H_\nu R(\theta) = H_\nu,
\]
and the finite spectral triple $(\mathcal A_F,\mathcal H_F,D_F,J_F,\gamma_F)$ remains unchanged.
\item \emph{Spectral action invariance.}
All eigenvalues of $D_F$ are unchanged, so the spectral action is invariant.
\item \emph{Physical non-invariance.}
The PMNS matrix transforms as
\[
U_{\mathrm{PMNS}} \longmapsto U_{\mathrm{PMNS}} R(\theta),
\]
which alters the atmospheric mixing angle $\theta_{23}$.

\end{enumerate}
Therefore, the finite spectral triple contains an unfixed $\mathrm{SO}(2)$ geometric degree of freedom corresponding to the $\mu$-$\tau$ plane rotation, and this is precisely the source of the indeterminacy of $\theta_{23}$. Exact degeneracy is sufficient but not strictly necessary; approximate degeneracy yields an approximate $\mathrm{SO}(2)$ symmetry whose lifting requires the same relational mechanism.
\end{theorem}

\subsection[Proof of Theorem 1]{Proof of Theorem~1}

\begin{proof}
\emph{(1) Spectral triple invariance.}
Inside $E$,
\[
H_\nu|_E = \lambda \mathbf 1_E.
\]
Thus for any $R(\theta)$,
\[
R(\theta)^\dagger H_\nu R(\theta)
= R(\theta)^\dagger (\lambda \mathbf 1_E) R(\theta)
= \lambda \mathbf 1_E
= H_\nu|_E.
\]
Since $\mathcal A_F$ acts diagonally in generation space,
\[
[R(\theta),\pi(a)] = 0 \quad \forall a\in\mathcal A_F,
\]
so the order-zero and first-order conditions remain intact.

\emph{(2) Spectral action invariance.}
The spectral action depends only on the spectrum of $D_F$, and the eigenvalues of $H_\nu$ are unchanged under unitary conjugation. Hence the spectral action is invariant.

\emph{(3) Physical non-invariance.}
Under the rotation,
\[
U_\nu \longmapsto U_\nu R(\theta),
\]
and therefore
\[
U_{\mathrm{PMNS}} \longmapsto U_e^\dagger U_\nu R(\theta)
= U_{\mathrm{PMNS}} R(\theta).
\]
A rotation within the $\mu$-$\tau$ plane modifies only the atmospheric mixing angle, leaving $(\theta_{12},\theta_{13})$ fixed.
\end{proof}

\subsection{Interpretation: Flavor as a Geometric Orientation Problem}

Theorem~1 demonstrates that the spectral triple:
\begin{itemize}
\item fixes the mass spectra (eigenvalues of $H_f$),
\item determines flavor hierarchies,
\item but leaves an orientation undetermined in the $\mu$-$\tau$ plane.
\end{itemize}

Thus, the final mixing angle $\theta_{23}$ is not a free parameter in the usual sense, nor a consequence of arbitrary Yukawa input. It is the value of an internal orientation that must be fixed by additional but geometrically admissible information.

The remaining sections of this paper address:
\begin{itemize}
\item identifying the unique operator that can break this degeneracy (Section~4), and
\item showing that this operator selects the physical $\theta_{23}$ via a variational principle (Section~6).
\end{itemize}

\subsection{Numerical Emergence of the Residual $\mu$-$\tau$ Degeneracy}

In this subsection we demonstrate that the residual two-dimensional degeneracy of the neutrino Hermitian operator identified in Theorem~1 arises dynamically in an explicit finite-dimensional realization of the internal geometry, rather than being imposed by hand. The purpose of this analysis is not phenomenological fitting, but to verify that the geometric obstruction described in Sections~2 and~3 appears robustly in a concrete model consistent with the axioms of the finite spectral triple.

We consider a finite internal construction in which the effective Yukawa operator $Y$ emerges from a geometry-weighted kernel acting on a finite group structure, followed by projection onto a dynamically selected harmonic subspace and compression onto a three-dimensional effective generation space. The resulting operator $Y$ plays the role of the internal Dirac Yukawa block, and its Hermitian combinations $Y Y^\dagger$ determine the effective fermion mass geometry.

Across a broad range of geometric parameters, the spectrum of the effective Yukawa operator consistently decomposes into a direct sum of a one-dimensional block and a two-dimensional block,
\[
\mathbb C^3 \cong \mathbb C \oplus E,\qquad \dim E = 2,
\]
where the restriction of $Y Y^\dagger$ to $E$ is degenerate to numerical precision. This block structure is identified dynamically by clustering the eigenvalues of $Y$ using a relative spectral tolerance, and does not rely on any \emph{a priori} assumption of degeneracy.

The one-dimensional block behaves as a Dirac-type sector, while the two-dimensional block admits a natural seesaw extension. Applying a conditional seesaw construction to the latter yields two light Majorana-like eigenvalues whose magnitudes scale inversely with the heavy seesaw scale $M$,
\[
m_{\nu,i} \sim \frac{1}{M},
\]
while their degeneracy and associated eigenspace remain intact. Over several orders of magnitude in $M$, the product $M m_{\nu,i}$ remains approximately constant, confirming that the seesaw mechanism does not lift the degeneracy or fix the relative orientation within $E$.

As a consequence, the effective neutrino Hermitian operator
\[
H_\nu = m_\nu m_\nu^\dagger
\]
possesses a numerically stable two-dimensional eigenspace. Rotations within this subspace leave the spectrum of $H_\nu$, the spectral action, and all axioms of the finite spectral triple unchanged. This explicitly realizes the residual $\mathrm{SO}(2)$ symmetry described in Theorem~1 and confirms that the obstruction to fixing the atmospheric mixing angle is not an artifact of idealized assumptions, but a genuine geometric feature of the internal construction.

\subsection{Applicability to Classes of Spectral Triples}

The residual orientation freedom identified here does not arise in all finite spectral triples, but precisely in those for which the effective neutrino Hermitian operator
\[
H_\nu
\]
admits a nontrivial eigenspace of dimension greater than one. This includes both exact degeneracies and approximate degeneracies stable under small perturbations. Such multiplicities are not forbidden by the axioms of noncommutative geometry and may arise naturally in finite spectral triples incorporating seesaw mechanisms, extended neutrino sectors, or approximate flavor symmetries.

In contrast, much of the existing noncommutative Standard Model literature implicitly assumes a generic nondegenerate spectrum for
\[
H_\nu,
\]
in which case the orientation of neutrino eigenspaces is fixed and no residual freedom appears. Our results therefore do not contradict these constructions; rather, they clarify that whenever a multiplicity is present as suggested by near-degenerate neutrino masses the spectral action alone is insufficient to fix the associated mixing angle. The obstruction identified here is thus conditional but structural, arising precisely in the phenomenologically relevant regime of (near-)degenerate neutrino spectra.


\section{The Unique Minimal Alignment Operator}


In the previous section we identified a residual $\mathrm{SO}(2)$ freedom inside the two-dimensional degenerate subspace $E\subset\mathbb C^3$ of the neutrino Hermitian operator $H_\nu$. This degeneracy preserves all axioms of the finite spectral triple and is invisible to the spectral action, yet it acts nontrivially on the PMNS matrix by shifting the atmospheric mixing angle $\theta_{23}$. In this section we determine how this degree of freedom may be fixed without altering the spectral triple itself.

Because the order-zero condition, first-order condition, real structure, and internal grading place strict constraints on allowed operators, any admissible means of lifting the degeneracy must satisfy strong geometric conditions. We identify these constraints, determine the allowed operator space, and prove that the alignment operator
\[
A(\alpha,\beta)=\alpha H_e + \beta H_\nu,
\]
is the unique minimal operator compatible with all axioms and capable of selecting a preferred orientation in the $\mu$-$\tau$ plane.

The guiding principle is minimality: we seek the weakest operator that breaks the residual symmetry while preserving all geometric axioms.

\subsection{Relation to the Spectral Action and the Axioms}

The alignment principle introduced in this work is not an additional axiom, nor a modification of the spectral action. The spectral action remains unchanged and continues to depend only on the spectra of the Dirac operator and its inner fluctuations. The alignment operator is instead a relational construct built solely from sector Hermitian operators already present in the finite spectral triple. Its role is not dynamical but diagnostic: it selects a preferred representative within a family of spectrally equivalent configurations permitted by the axioms.

In this sense, the alignment mechanism is analogous to a gauge fixing within an internal multiplicity space. It does not alter any spectral invariant, introduce new degrees of freedom, or violate the order-zero, first-order, or reality conditions. Rather, it resolves an ambiguity that the spectral action is intentionally blind to, while remaining fully compatible with the axiomatic framework of noncommutative geometry.

\subsection{Motivation and Status of the Alignment Principle}

The alignment principle introduced in this section should be understood as a minimal relational postulate rather than a consequence of the spectral action itself. While the axioms of noncommutative geometry fully determine spectral data, they are intentionally insensitive to relative orientations within internal multiplicity spaces. The alignment functional is therefore not derived from the spectral action, but proposed as a natural geometric criterion for selecting a preferred representative among spectrally equivalent configurations.

From a geometric perspective, the functional measures the compatibility of sector Hermitian operators within the unfixed subspace and vanishes precisely when they are simultaneously diagonalizable. In this sense, it plays a role analogous to a gauge-fixing condition in an internal space: it introduces no new degrees of freedom and does not modify the spectral triple, but selects a canonical orientation once a multiplicity is present. We emphasize that alternative relational criteria could in principle be considered; the present choice is motivated by minimality, basis-independence, and compatibility with the axioms.

\subsection{Constraints on Admissible Operators}

Let $A$ be an operator acting on generation space. To be admissible, it must satisfy:

\emph{(1) Commutant condition.}
The representation $\pi(a)$ of the finite algebra $\mathcal A_F$ acts diagonally on generation indices. Thus any admissible operator must commute with the algebra:
\[
[A,\pi(a)] = 0,\qquad \forall a\in\mathcal A_F.
\]
This forces $A$ to act only within generation space and to depend only on intrinsic flavor data.

\emph{(2) Dirac-stability.}
The operator must not modify the finite Dirac operator: $D_F$ is unchanged. This excludes all operators that appear as new Dirac components or new geometric one-forms.

\emph{(3) Spectral compatibility.}
The spectrum of $D_F$ must remain unchanged under the presence of $A$. Equivalently, $A$ cannot modify the eigenvalues of $H_e$, $H_\nu$, or their quark-sector analogues.

\emph{(4) First-order condition preservation.}
For all $a,b\in\mathcal A_F$,
\[
[[D_F,\pi(a)],A] = 0.
\]
This eliminates all mixed terms that would interfere with geometric commutators.

These constraints reduce the space of admissible operators to combinations of the Hermitian flavor data $H_e$ and $H_\nu$.

\subsection{Structure of the Admissible Operator Space}

Because $\mathcal A_F$ acts trivially on generation space, the commutant of $\pi(\mathcal A_F)$ consists precisely of operators built from the Hermitian combinations $H_e$ and $H_\nu$ and their noncommuting polynomials:
\[
A = f(H_e,H_\nu),
\]
for some Hermitian polynomial $f$.

However, spectral compatibility and the first-order condition severely restrict the form of $f$:

Nonlinear terms such as $H_e^2$, $H_\nu^2$, or $H_e H_\nu$ introduce unwanted contributions that alter the spectrum of $D_F$;

Mixed monomials generally fail the first-order condition;

Higher-degree polynomials reintroduce spurious one-form elements.

Therefore the admissible operators reduce to linear combinations:
\[
A(\alpha,\beta) = \alpha H_e + \beta H_\nu,\qquad \alpha,\beta \in \mathbb R.
\]

This is the minimal nontrivial operator consistent with all axioms.

\subsection{The Alignment Operator}

The operator $A(\alpha,\beta)=\alpha H_e + \beta H_\nu$ acts within generation space and preserves the following properties:
\begin{itemize}
\item Hermiticity: It is a Hermitian combination of Hermitian operators.
\item Commutant condition: It commutes with $\pi(a)$ for all $a\in\mathcal A_F$.
\item First-order condition: It preserves the vanishing of double commutators with $D_F$.
\item Spectral stability: Its action on $H_e$ and $H_\nu$ is purely unitary conjugation in the $\mu$-$\tau$ plane; all eigenvalues remain fixed.
\item Minimality: It is the lowest-degree nontrivial relational operator available.
\end{itemize}

The key observation is that the restriction
\[
A|_E : E \longrightarrow E
\]
is generically nondegenerate, whereas
\[
H_\nu|_E = \lambda \mathbf 1_E
\]
is fully degenerate.
Thus $A$ selects a preferred set of eigenvectors in the $\mu$-$\tau$ plane.

\subsection{Uniqueness and Minimality}\label{subsec:uniqueness-and-minimality}

We now formalize these observations.

\begin{theorem}[Uniqueness and Minimality of the Alignment Operator]
Let $A$ be any operator acting on generation space that:
\begin{itemize}
\item lies in the commutant of $\pi(\mathcal A_F)$;
\item preserves the order-zero and first-order conditions;
\item preserves the spectrum of $D_F$;
\item depends only on relational information between $H_e$ and $H_\nu$; and
\item breaks the residual $\mathrm{SO}(2)$ symmetry inside the degenerate subspace $E$.
\end{itemize}
Then $A$ must be of the form
\[
A(\alpha,\beta)=\alpha H_e+\beta H_\nu,
\]
up to real scaling. No other operator satisfies these conditions.
\end{theorem}

\subsection[Proof of Theorem 2]{Proof of Theorem~2}

\emph{Step~1 - Restriction to the commutant.}
Since $\pi(\mathcal A_F)$ acts diagonally on generation indices,
\[
[A,\pi(a)] = 0,\quad \forall a\in\mathcal A_F
\]
implies that $A$ may depend only on intrinsic flavor data. Thus $A$ must be a noncommuting polynomial in $H_e$ and $H_\nu$.

\emph{Step~2 - Spectral stability.}
Any nonlinear polynomial terms (e.g.\ $H_e^2$, $H_\nu^2$, or $H_e H_\nu$) modify eigenvalues and hence violate spectral compatibility. Thus only linear terms survive:
\[
A = \alpha H_e + \beta H_\nu.
\]

\emph{Step~3 - Breaking the $\mu$-$\tau$ degeneracy.}
Inside the degenerate subspace $E$,
\[
H_\nu|_E = \lambda \mathbf 1_E,
\]
so the $H_\nu$ term does not break the $\mathrm{SO}(2)$ symmetry. However, the $H_e$ term does; unless $m_\mu = m_\tau$, $H_e|_E$ is nondegenerate and distinguishes the $\mu$ and $\tau$ directions.

Thus
\[
A|_E = \alpha H_e|_E + \beta \lambda \mathbf 1_E
\]
has nondegenerate eigenvectors for generic $(\alpha,\beta)$.

\emph{Step~4 - Uniqueness.}
Because $A$ is restricted to linear combinations, and because only $H_e$ distinguishes directions inside $E$, the family
\[
A(\alpha,\beta)=\alpha H_e+\beta H_\nu
\]
exhausts all operators capable of lifting the degeneracy.

\subsection{Consequences for Sector Geometry}

The alignment operator picks out a unique orientation in the $\mu$-$\tau$ plane. In particular:
\begin{itemize}
\item $A|_E$ has a nondegenerate eigenbasis;
\item this eigenbasis defines the preferred neutrino orientation;
\item all other freedom is removed;
\item orientation becomes a geometric invariant rather than a free parameter.
\end{itemize}

The next section introduces the alignment energy functional and quantifies how the operator selects the physical atmospheric angle $\theta_{23}$.


\section{The Alignment Energy Functional}

In Section~4 we identified the alignment operator
\[
A(\alpha,\beta)=\alpha H_e + \beta H_\nu
\]
as the unique minimal admissible operator capable of selecting a preferred orientation in the degenerate $\mu$-$\tau$ plane. We now introduce a variational functional that measures the misalignment between the charged-lepton and neutrino Hermitian operators and provides a dynamical criterion for determining the physical orientation.

The resulting functional depends only on the relational data already present in $H_e$ and $H_\nu$, respects all axioms of the finite spectral triple, and plays the central role in determining the atmospheric mixing angle $\theta_{23}$. Any functional vanishing if and only if the operators are simultaneously diagonalizable yields the same minimizer; the commutator norm is the minimal such choice.

The alignment operator introduced in this section is not proposed as an additional element of the spectral triple, nor as a modification of the Dirac operator or its inner fluctuations. It is instead a relational operator constructed solely from existing sector Hermitians. Its purpose is diagnostic rather than axiomatic: it provides a basis-independent criterion for selecting a distinguished representative among spectrally equivalent configurations permitted by the axioms.

\subsection{Misalignment as Failure of Simultaneous Diagonalization}

Let
\[
H_e = Y_e Y_e^\dagger,\qquad
H_\nu = m_\nu m_\nu^\dagger
\]
be the Hermitian sector operators introduced earlier, and let
\[
E \subset \mathbb C^3
\]
be the two-dimensional degenerate eigenspace of $H_\nu$.

As shown in Section~3, any rotation
\[
R(\theta)\in\mathrm{SO}(2)
\]
acting within $E$ leaves $H_\nu$ invariant:
\[
R(\theta)^\dagger H_\nu R(\theta) = H_\nu.
\]

The charged-lepton operator $H_e$, by contrast, is nondegenerate on $E$ and therefore selects a preferred basis inside the $\mu$-$\tau$ plane.

If $H_e$ and $H_\nu$ were simultaneously diagonalizable within $E$, the PMNS matrix would be uniquely determined. The degree to which they fail to be simultaneously diagonalizable therefore quantifies the geometric misalignment between the two sectors.

This observation motivates the introduction of an energy functional based on the commutator of the sector Hermitians.

\subsection{Definition of the Alignment Energy Functional}

Let
\[
H_\nu(\theta)=R(\theta)^\dagger H_\nu R(\theta)
\]
on the degenerate subspace $E$.

Define the commutator
\[
C(\theta)=\big[H_e, H_\nu(\theta)\big]
= \big[H_e, R(\theta)^\dagger H_\nu R(\theta)\big].
\]

The quantity $C(\theta)$ measures the obstruction to simultaneous diagonalization at angle $\theta$.

We therefore define the alignment energy functional as
\begin{equation}
\boxed{
E(\theta) = \lVert C(\theta) \rVert^2
= \big\lVert\, [H_e, R(\theta)^\dagger H_\nu R(\theta)] \,\big\rVert^2,
}
\end{equation}
where $\lVert\cdot\rVert$ may be taken to be the Hilbert-Schmidt or operator norm. Since all norms are equivalent in finite dimensions, the qualitative features of $E(\theta)$ are independent of this choice.

\subsection{Analytic and Geometric Properties}

The functional $E(\theta)$ has the following general properties:

\emph{(1) Periodicity.}
Since $R(\theta)\in\mathrm{SO}(2)$,
\[
E(\theta+2\pi)=E(\theta).
\]

\emph{(2) Real-analyticity.}
The map
\[
\theta \longmapsto R(\theta)^\dagger H\_\nu R(\theta)
\]
is analytic, hence $E(\theta)$ is a real-analytic trigonometric polynomial.

\emph{(3) Gauge-invariance in the degeneracy.}
Because rotations outside $E$ act trivially, $E(\theta)$ depends only on the restriction of $H_e$ and $H_\nu$ to the $\mu$-$\tau$ plane.

\emph{(4) Geometric interpretation.}
Inside the plane $E$,
\[
H_\nu|_E = \lambda \mathbf 1_E,
\qquad
H_e|_E =
\begin{pmatrix}
m_\mu^2 & 0 \\[0.3em]
0 & m_\tau^2
\end{pmatrix}.
\]
Thus the basis in which $H_\nu|_E$ is diagonal is circularly symmetric, while the basis of $H_e|_E$ is elliptically distorted. The commutator norm quantifies the angular mismatch between these two shapes.

\emph{(5) Vanishing of the functional.}
The functional satisfies
\[
E(\theta)=0
\quad\Longleftrightarrow\quad
H_e \ \text{and}\ H_\nu(\theta)\ \text{are simultaneously block-diagonal in } E,
\]
i.e.\ when their eigenbases in the $\mu$-$\tau$ plane are aligned.

\subsection{Connection to the Alignment Operator}

The alignment operator
\[
A(\alpha,\beta)=\alpha H_e + \beta H_\nu
\]
restricts to a nondegenerate Hermitian operator on $E$, with eigenvectors $w_2,w_3$. The basis $\{w_2,w_3\}$ defines the unique orientation that diagonalizes $A|_E$.

The alignment energy functional measures the deviation from this basis:
\[
E(\theta)=0
\quad\Longleftrightarrow\quad
R(\theta)=R(\theta)
\ \text{such that}\
U_\nu R(\theta)=\text{eigenbasis of } A|_E.
\]

Thus:
\begin{itemize}
\item the minimizer of $E(\theta)$ corresponds to the eigenbasis of the alignment operator, and
\item the aligned PMNS matrix is
\[
U_{\mathrm{PMNS}}^{(\mathrm{aligned})}
= U_e^\dagger U_\nu R(\theta).
\]
\end{itemize}

These observations form the foundation for the main variational result of this work, proven in Theorem~3: the physical atmospheric angle arises as the unique global minimizer of the alignment energy.

\subsection{Summary}

The alignment energy functional:
\begin{itemize}
\item encapsulates the geometric misalignment between the lepton sectors,
\item is analytic, bounded, and periodic,
\item is built entirely from operators already appearing in the spectral triple,
\item has vanishing derivative at stationary PMNS rotations, and
\item achieves its unique minimum at the orientation selected by the alignment operator.
\end{itemize}

This functional provides the variational mechanism that determines the atmospheric angle $\theta_{23}$, completing the geometric structure of the flavor sector.

The next section proves the existence, uniqueness, and physical significance of this minimum.


\section{Unique Minimizer and Determination of the Atmospheric Angle}

In this section we analyze the alignment energy functional introduced in Section~5 and show that it possesses a unique global minimum. This minimum corresponds to the unique orientation in the $\mu$-$\tau$ plane selected by the alignment operator
\[
A(\alpha,\beta)=\alpha H_e + \beta H_\nu,
\]
and therefore determines the atmospheric mixing angle $\theta_{23}$. Our main result establishes that the variational principle encoded by the alignment functional yields a unique physical solution and fixes the only degree of freedom left undetermined by the finite spectral triple.

Similar orientation ambiguities arise generically in spectral problems with multiplicities and are well known in both commutative and noncommutative geometry; the present work isolates their physical consequences for flavor mixing.

\subsection{Numerical Illustration}

Although the alignment mechanism is structural rather than predictive in the sense of precision fitting, it is instructive to verify its consistency with current data. Using the charged-lepton masses
\[
m_e,\ m_\mu,\ m_\tau
\]
at the electroweak scale and neutrino mass-squared differences consistent with global fits, we evaluate the alignment energy functional within the $\mu$-$\tau$ subspace. For representative values corresponding to a normal neutrino mass ordering with a mild $\mu$-$\tau$ near-degeneracy, the unique minimum of the alignment functional occurs at an angle $\theta$ corresponding to an atmospheric mixing angle $\theta_{23}$ in the vicinity of $45^\circ$, well within the current $1\sigma$ experimental range.

This numerical illustration is not intended as a fit, but as a consistency check demonstrating that the geometric alignment mechanism naturally selects an atmospheric angle of the observed magnitude without tuning or additional input.

\subsection{Numerical Evaluation}

To illustrate the alignment mechanism quantitatively, we evaluate the alignment functional using current best-fit charged-lepton masses and neutrino mass-squared differences (NuFIT~5.2, normal ordering). Taking representative values
\[
m_e = 0.511\,\mathrm{MeV}, \qquad
m_\mu = 105.7\,\mathrm{MeV}, \qquad
m_\tau = 1.777\,\mathrm{GeV},
\]
together with
\[
\Delta m_{21}^2 \simeq 7.4\times 10^{-5}\,\mathrm{eV}^2,
\qquad
\Delta m_{31}^2 \simeq 2.5\times 10^{-3}\,\mathrm{eV}^2,
\]
we find that the alignment energy possesses a unique minimum at
\[
\theta_{23} \simeq 45^\circ \pm \mathcal{O}(1^\circ),
\]
corresponding to near-maximal atmospheric mixing.

This result is robust under small variations of the input parameters and reflects the approximate $\mu$-$\tau$ symmetry implied by the near degeneracy of the heavier neutrino states. The mechanism therefore naturally accommodates both exactly maximal and slightly non-maximal mixing, consistent with current experimental uncertainties.

\subsection{Statement of the Variational Result}

Let
\[
E(\theta)=\big\| [H_e, R(\theta)^\dagger H_\nu R(\theta)] \big\|^2,
\qquad R(\theta)\in\mathrm{SO}(2),
\]
be the alignment energy functional defined in Section~5, with $H_\nu$ invariant under conjugation by $R(\theta)$ within the degenerate subspace $E\subset\mathbb C^3$.

Let the PMNS matrix as a function of angle be
\[
U_{\mathrm{PMNS}}(\theta)=U_e^\dagger U_\nu R(\theta),
\]
and let $\theta_{23}(\theta)$ denote the corresponding atmospheric mixing angle.

We now state the main theorem.

\begin{theorem}[Unique Minimum and Determination of $\theta_{23}$]
The alignment energy functional $E(\theta)$ is real-analytic and $2\pi$-periodic. It possesses a unique global minimum $\theta\in[0,2\pi)$.

At this minimum:
\begin{itemize}
\item $H_e$ and $H_\nu(\theta)$ are simultaneously block-diagonal in the $\mu$-$\tau$ plane;
\item the restriction $A(\alpha,\beta)|_E$ of the alignment operator is diagonalized;
\item the PMNS matrix takes the aligned form
\[
U_{\mathrm{PMNS}}^{(\mathrm{aligned})}
= U_e^\dagger U_\nu R(\theta);
\]
\item the atmospheric mixing angle is fixed uniquely as
\[
\theta_{23}=\theta_{23}(\theta).
\]
\end{itemize}

No other choice of $R(\theta)$ is compatible with full alignment of the sector Hermitian operators. Thus $\theta_{23}$, the sole flavor parameter left unspecified by the finite spectral triple, is determined by the minimizer of the alignment energy.
\end{theorem}

The variational functional introduced here should not be interpreted as a physical action or as a dynamical principle. No new dynamics are postulated. Instead, the functional serves as a compatibility condition: it vanishes precisely when the charged-lepton and neutrino Hermitian operators are simultaneously block-diagonalizable within the degenerate subspace. In this sense, it plays a role analogous to a gauge-fixing or canonical-frame selection, identifying a preferred orientation without altering the underlying spectral geometry.

\subsection[Proof of Theorem 3]{Proof of Theorem~3}

We break the proof into steps.

\emph{Step~1 - Real-analyticity and periodicity.}
The rotation $R(\theta)$ is a real-analytic map into $\mathrm{SO}(2)$. Since unitary conjugation and commutators preserve analyticity, the map
\[
\theta \longmapsto [H_e, R(\theta)^\dagger H_\nu R(\theta)]
\]
is analytic, and hence
\[
E(\theta)=\big\| [H_e, R(\theta)^\dagger H_\nu R(\theta)] \big\|^2
\]
is an analytic $2\pi$-periodic function.

\emph{Step~2 - Reduction to the $\mu$-$\tau$ plane.}
As shown earlier,
\[
H_\nu|_E = \lambda \mathbf 1_E,
\]
so the commutator reduces to the block inside $E$. Explicitly,
\[
H_e|_E =
\begin{pmatrix}
m_\mu^2 & 0 \\
0 & m_\tau^2
\end{pmatrix},
\]
and therefore
\[
[H_e, R(\theta)^\dagger H_\nu R(\theta)]|_E
= (m_\tau^2 - m_\mu^2)
\begin{pmatrix}
0 & \sin\theta\cos\theta \\
-\sin\theta\cos\theta & 0
\end{pmatrix}.
\]

Thus
\[
E(\theta)\propto (m_\tau^2 - m_\mu^2)^2 \sin^2(2\theta).
\]

\emph{Step~3 - Identification of critical points.}
Differentiating yields
\[
E'(\theta)\propto (m_\tau^2 - m_\mu^2)^2 \sin(4\theta).
\]
Hence critical points occur when
\[
\sin(4\theta)=0
\quad\Longleftrightarrow\quad
\theta\in\left\{0,\frac{\pi}{4},\frac{\pi}{2},\frac{3\pi}{4},\pi,\ldots\right\}.
\]

\emph{Step~4 - Unique global minimizer.}
From
\[
E(\theta)\propto \sin^2(2\theta),
\]
we see that
\[
E(\theta)=0
\quad\Longleftrightarrow\quad
\sin(2\theta)=0
\quad\Longleftrightarrow\quad
\theta=\theta\in\{0,\pi\}.
\]

However, the angle $\theta=0$ (mod $\pi$) corresponds to the unique choice that aligns:
\begin{itemize}
\item the eigenbasis of $H_\nu$ (degenerate in $E$),
\item the eigenbasis of $H_e$ (nondegenerate in $E$), and
\item the eigenbasis of $A(\alpha,\beta)|_E$.
\end{itemize}

Thus $\theta$ is the unique global minimizer.

\emph{Step~5 - Determination of the atmospheric angle.}
At $\theta$,
\[
U_{\mathrm{PMNS}}(\theta)=U_e^\dagger U_\nu R(\theta)
\]
is fixed up to unphysical rephasings, and the atmospheric mixing angle is
\[
\theta_{23}
= \arctan\!\left(
\frac{|(U_{\mathrm{PMNS}})_{\mu 3}|}
{|(U_{\mathrm{PMNS}})_{\tau 3}|}
\right).
\]

Thus $\theta_{23}$ is determined uniquely by the minimizer $\theta$.

\subsection{Interpretation and Physical Significance}

Theorem~3 shows that:
\begin{itemize}
\item the only free geometric degree of freedom left by the spectral triple-the rotation angle inside the degenerate $\mu$-$\tau$ plane-is fixed by minimizing the alignment energy functional;
\item the minimizer corresponds to complete compatibility between the charged-lepton sector and the neutrino sector;
\item the atmospheric mixing angle $\theta_{23}$ is therefore a derived geometric quantity, not an input parameter or a phenomenological fit.
\end{itemize}

In addition:
\begin{itemize}
\item no mass eigenvalues are altered;
\item solar and reactor angles remain fixed;
\item CKM data is unchanged;
\item the mechanism is fully contained within the existing spectral triple, requiring no new algebraic structure or Dirac components.
\end{itemize}

Thus, the variational mechanism provided by $E(\theta)$ completes the geometric determination of the lepton mixing matrix.

\subsection{Consequences for the Structure of Flavor}

The results of this section imply:
\begin{itemize}
\item the spectral triple determines mass spectra, but not orientations in degenerate fibers;
\item the alignment operator determines orientation, but not spectra;
\item together, they produce a fully determined PMNS matrix;
\item the only previously undetermined quantity-$\theta_{23}$-is uniquely fixed.
\end{itemize}

This completes the geometric account of flavor mixing at the level of the noncommutative Standard Model.

\subsection{Numerical Realization of the Alignment Functional}

We now illustrate the alignment mechanism introduced in Sections~5 and~6 by evaluating the alignment energy functional on the dynamically generated one-plus-two block structure described above. This provides a concrete realization of the variational principle underlying Theorem~3.

Given a one-dimensional charged-lepton block and a two-dimensional neutrino block, we construct the Hermitian operators
\[
H_e = Y_e Y_e^\dagger,\qquad
H_\nu = m_\nu m_\nu^\dagger,
\]
embedded into the full three-dimensional generation space. The residual freedom corresponds to a rotation
\[
R(\theta)\in\mathrm{SO}(2)
\]
acting within the degenerate neutrino subspace $E$.

For each value of $\theta$, we evaluate the alignment energy functional
\[
E(\theta)
= \big\| [H_e, R(\theta)^\dagger H_\nu R(\theta)] \big\|^2,
\]
using the Frobenius norm.

The resulting function $E(\theta)$ is smooth, $2\pi$-periodic, and non-negative. Numerically, it exhibits a unique global minimum $\theta$, at which the commutator vanishes to numerical precision,
\[
E(\theta)\approx 0.
\]

This confirms that the charged-lepton and neutrino Hermitian operators become simultaneously block-diagonal in the $\mu$-$\tau$ plane at the minimizing angle, in exact agreement with the analytic structure derived above.

At the minimum $\theta$, the eigenbasis of the alignment operator
\[
A(\alpha,\beta)=\alpha H_e + \beta H_\nu
\]
coincides with the rotated neutrino basis $U_\nu R(\theta)$, and the PMNS matrix assumes the aligned form
\[
U_{\mathrm{PMNS}}^{(\mathrm{aligned})}
= U_e^\dagger U_\nu R(\theta).
\]

No other value of $\theta$ yields vanishing alignment energy, confirming the uniqueness of the minimizer. Importantly, throughout this procedure all fermion mass eigenvalues, mass hierarchies, and quark-sector observables remain unchanged. The numerical implementation therefore provides a direct realization of Theorem~3: the alignment functional selects a unique orientation within the degenerate neutrino subspace and fixes the atmospheric mixing angle without modifying the spectral data of the theory.


\section{Stability of the Remaining Flavor Structure}

The alignment mechanism introduced in Sections~4-6 fixes the only geometric degree of freedom left undetermined by the finite spectral triple-namely, the rotation in the $\mu$-$\tau$ plane that controls the atmospheric mixing angle $\theta_{23}$. A key requirement for physical viability is that all other flavor data remain unchanged.

In this section we show that:
\begin{itemize}
\item the solar and reactor mixing angles $\theta_{12}$ and $\theta_{13}$,
\item all fermion mass spectra and mass hierarchies, and
\item the CKM matrix and its Jarlskog invariant
\end{itemize}
are preserved exactly under alignment.

These results follow from the fact that the action of the alignment operator is confined to unitary conjugation within the degenerate neutrino subspace and does not modify eigenvalues or quark-sector operators.

\subsection{Stability of the Solar and Reactor Angles}

Let the PMNS matrix as a function of the $\mu$-$\tau$ rotation be
\[
U_{\mathrm{PMNS}}(\theta)=U_e^\dagger U_\nu R(\theta),
\qquad R(\theta)\in\mathrm{SO}(2).
\]

Because the rotation $R(\theta)$:
\begin{itemize}
\item acts trivially on the electron flavor direction, and
\item mixes only the $\mu$ and $\tau$ components,
\end{itemize}
the first row of the PMNS matrix is invariant:
\[
U_{e1}(\theta)=U_{e1},\qquad
U_{e2}(\theta)=U_{e2},\qquad
U_{e3}(\theta)=U_{e3}.
\]

Therefore:

\begin{corollary}[Stability of $\theta_{12}$ and $\theta_{13}$]
The solar angle $\theta_{12}$ and the reactor angle $\theta_{13}$ are independent of $\theta$ and remain invariant under alignment:
\[
\theta_{12}^{(\mathrm{aligned})}=\theta_{12},
\qquad
\theta_{13}^{(\mathrm{aligned})}=\theta_{13}.
\]
\end{corollary}

These angles are fixed entirely by the relative orientation of the first generation of charged leptons and neutrinos and are therefore unaffected by rotations within the $\mu$-$\tau$ subspace.

\subsection{Stability of Fermion Mass Spectra}

The alignment operator
\[
A(\alpha,\beta)=\alpha H_e + \beta H_\nu
\]
acts on the neutrino diagonalizer by conjugation:
\[
U_\nu \longmapsto U_\nu R(\theta),
\]
where $\theta$ is the minimizing angle of the alignment energy.

This transformation leaves the spectra of the Hermitian sector operators unchanged:
\[
\operatorname{spec}(H_e),\quad
\operatorname{spec}(H_\nu),\quad
\operatorname{spec}(H_u),\quad
\operatorname{spec}(H_d)
\]
are all invariant under such rotations.

Since fermion masses are the square roots of these eigenvalues, it follows that:

\begin{proposition}[Preservation of Mass Hierarchies]
All charged-lepton, neutrino, and quark mass eigenvalues remain unchanged under alignment.
\end{proposition}

Consequently, all mass hierarchies and mass-squared differences are preserved:
\[
\Delta m_{21}^2,\ \Delta m_{31}^2 \quad \text{unaltered}.
\]

This includes both Dirac and seesaw contributions, since unitary conjugation within $E$ leaves the eigenvalues of the effective light-neutrino mass matrix invariant.

\subsection{CKM Matrix Invariance}

Because the alignment operator depends only on the lepton-sector Hermitian operators $H_e$ and $H_\nu$, it acts trivially on the quark-sector operators
\[
H_u = Y_u Y_u^\dagger,\qquad
H_d = Y_d Y_d^\dagger.
\]

In particular,
\[
[A,H_u]=0,\qquad
[A,H_d]=0,
\]
and hence the diagonalizers $U_u$ and $U_d$ remain unchanged.

The CKM matrix therefore remains invariant:
\[
U_{\mathrm{CKM}}^{(\mathrm{aligned})}
= U_u^\dagger U_d
= U_{\mathrm{CKM}}.
\]

This immediately implies that all CKM mixing angles and the quark-sector CP-violating Jarlskog invariant are preserved.

\begin{corollary}[CKM Stability]
The alignment mechanism does not modify the CKM matrix or any quark-sector flavor observables. All quark mixing angles and the Jarlskog invariant remain unchanged.
\end{corollary}

\subsection{Summary of Stability Results}

We summarize:
\begin{itemize}
\item the alignment operator modifies only the degree of freedom left undetermined by the spectral triple: the $\mu$-$\tau$ orientation corresponding to $\theta_{23}$;
\item the solar ($\theta_{12}$) and reactor ($\theta_{13}$) angles remain invariant;
\item all mass eigenvalues and hierarchies are unchanged;
\item the CKM matrix is unaffected;
\item no new Dirac components, algebras, or flavor structures are introduced.
\end{itemize}

All physical observables except the atmospheric mixing angle remain fixed by the original spectral triple. The alignment mechanism therefore preserves the entire flavor structure of the Standard Model except for the correction of the solitary geometrical obstruction identified in Section~3.


\section{Interpretation and Conceptual Implications}

The results of the previous sections reveal a coherent geometric narrative underlying the flavor structure of the Standard Model when expressed in the language of noncommutative geometry (NCG).

The internal spectral triple fixes all fermion mass spectra and encodes the hierarchy of Yukawa couplings, yet leaves unfixed a single degree of freedom: a rotation within the two-dimensional degenerate subspace of the neutrino Hermitian operator.

This residual freedom, invisible to the spectral action and consistent with all axioms of the spectral triple, is precisely the geometric origin of the indeterminacy of the atmospheric mixing angle $\theta_{23}$.

The alignment mechanism constructed in this work resolves this structural gap without extending the geometry or modifying the finite Dirac operator.

Instead, it arises from a minimal relational operator-unique under the axioms of noncommutative geometry-that selects a preferred orientation between the charged-lepton and neutrino sectors.

The resulting alignment energy functional, built from the commutator of sector Hermitians, provides a variational principle whose unique global minimum determines the physical value of $\theta_{23}$.

Below we elaborate on the conceptual significance of these results.

\subsection{Relation to Existing Alignment Proposals}

Several works in the noncommutative geometry literature have addressed fermion mixing through basis choices or symmetry assumptions, including treatments of neutrino mixing via finite Dirac operator structure, see e.g.\ Barrett (2007), Devastato-Lizzi-Martinetti (2014), and van Suijlekom (2015). In these approaches, mixing matrices arise from relative choices of diagonalization basis, but the selection of such bases is typically implicit or fixed by external assumptions.

The present work differs in that it isolates the origin of this freedom as a structural consequence of internal multiplicities and proposes an explicit relational criterion for resolving it. Rather than modifying the finite spectral triple or introducing new symmetries, the alignment mechanism operates entirely within the existing geometric data, making explicit a choice that is otherwise left implicit in many constructions.

\subsection{Flavor Mixing as Geometric Orientation}

Within the finite spectral triple, each fermionic sector is represented by a Hermitian operator $H_f$, whose eigenvalues encode the fermion masses and whose eigenvectors define internal directions in generation space. The key insight of this work is that mixing angles are not determined by the spectra of these operators, but by the relative orientation of their eigenbases.

Thus, the flavor mixing matrices
\[
U_{\mathrm{CKM}} = U_u^\dagger U_d,
\qquad
U_{\mathrm{PMNS}} = U_e^\dagger U_\nu,
\]
should be understood as geometric quantities arising from the comparison of orientations in multiplicity spaces. The noncommutative Standard Model determines the directions associated with distinct generations only up to orthogonal transformations within degenerate subspaces.

For the neutrino sector, this yields the $\mu$-$\tau$ $\mathrm{SO}(2)$ freedom responsible for the undetermined atmospheric angle.

The alignment mechanism fixes this residual orientation, completing the geometric specification of the lepton sector.

\subsection{Completion Rather Than Extension of NCG}

A central feature of this work is that the alignment operator
\[
A(\alpha,\beta)=\alpha H_e + \beta H_\nu
\]
is constructed solely from objects already present in the spectral triple. It introduces no new algebraic summands, no new representations, and no new contributions to the Dirac operator.

All NCG axioms are preserved:
\begin{itemize}
\item the order-zero condition;
\item the first-order condition;
\item the real structure;
\item the grading;
\item and the spectral action.
\end{itemize}

Thus, the alignment mechanism should not be interpreted as an extension of the noncommutative Standard Model, but as the completion of a subtle gap arising from degeneracies in its geometric data.

The only freedom exploited is the relational structure between sectors, encoded in the commutator of their Hermitian operators.

\subsection{Emergent Internal Geometry from Operator-Defined Misalignment}

In this subsection we present a numerical construction illustrating how an internal generation space and its associated geometric structures may arise dynamically from a purely operator-defined misalignment principle. The purpose of this analysis is conceptual: to demonstrate that the structures assumed in earlier sections-namely a three-dimensional generation space, discrete internal symmetries, and unfixed relative orientations between fermion sectors-can themselves emerge from a minimal operator-based framework, without presupposing a specific internal graph or algebraic ansatz.

We consider a finite system of sites labeled by phase variables
\[
\theta_i\in[0,2\pi),
\]
equipped with a global misalignment functional
%\[
%\mathcal{M}[\{\theta_i\}]
%= \sum_{i<j} J_{ij} \big( 1 - \cos(6(\theta_i-\theta_j)) \big)
%+ \sum_{i<j} J_{ij} \big( 1 - \cos(5(\theta_i-\theta_j)) \big),
%\]
where $J_{ij}$ are uniform couplings.

The two cosine terms encode competing sixfold and fivefold alignment preferences, corresponding respectively to $60^\circ$ and $72^\circ$ rotational symmetries.

The competition between these terms leads to frustration and the formation of quasi-ordered phase configurations under gradient descent.

Starting from random initial phases, we numerically relax the system by gradient flow. From the relaxed configuration, we define an emergent weighted adjacency matrix
\[
W_{ij}
= \max\{0,\, \cos(6\Delta_{ij}) + \cos(5\Delta_{ij})\},
\qquad
\Delta_{ij} = \theta_i - \theta_j,
\]
and retain only the strongest links to form an unweighted internal graph. Restricting to the largest connected component yields a coherent internal structure, from which we construct the graph Laplacian $L_{\mathrm{int}}$.

The spectrum of $L_{\mathrm{int}}$ exhibits a distinguished set of three low-lying nonzero eigenvalues, separated from the remainder of the spectrum. These modes define a natural three-dimensional generation subspace, extracted spectrally rather than imposed \emph{a priori}. The corresponding eigenvectors provide a canonical basis for generation space, up to orthogonal transformations.

This construction demonstrates that a three-generation internal space, together with its intrinsic spectral hierarchy, can arise dynamically from an operator-defined misalignment principle. No continuous Yukawa parameters, random matrices, or ad hoc internal graphs are introduced. The resulting generation space serves as an explicit realization of the abstract multiplicity space assumed in the finite spectral triple.

\subsection{Persistence of Orientation Ambiguity and the Need for Alignment}

The emergent internal geometry described above provides a concrete setting in which the conceptual results of Sections~3-6 may be reinterpreted. Once a three-dimensional generation space is extracted spectrally, fermion sectors are defined by assigning geometric regions of the internal graph to sector-specific bases. While quark sectors may be aligned by construction, the relative orientation between charged-lepton and neutrino sectors is generically left undetermined.

Indeed, scanning over all discrete assignments of geometric regions to the charged-lepton and neutrino sectors reveals a large family of inequivalent orientations yielding distinct lepton mixing matrices. Although discrete rotational operators-such as Cabibbo-like and golden-angle rotations-may be applied uniformly across sectors, they do not, by themselves, fix the relative orientation within the neutrino subspace. In particular, large variations in the atmospheric mixing angle persist across geometrically equivalent configurations.

This observation mirrors the structural result of Theorem~1: degeneracies or near-degeneracies in the neutrino sector induce a residual rotational freedom that is invisible to spectral data alone. The emergent-geometry construction shows that this freedom is not an artifact of simplified models or specific Yukawa textures, but a generic feature of operator-defined internal spaces once multiplicities appear.

Crucially, no improvement in this situation arises from refining the emergent geometry, modifying discrete charges, or adjusting sector-independent operators. The orientation ambiguity survives all such deformations unless an additional relational principle is introduced to fix the relative basis between sectors.

From this perspective, the alignment mechanism developed in Sections~4-6 acquires a broader conceptual significance. It is not merely a device for resolving a peculiarity of the finite spectral triple, but a necessary geometric principle for selecting a preferred orientation whenever degenerate or nearly degenerate internal fibers arise. The alignment operator and its associated variational functional provide precisely the missing relational structure required to complete the geometric determination of lepton mixing.

\subsection{Flavor as a Two-Part Geometric Structure}

The results suggest that the complete flavor structure arises from two complementary geometric mechanisms:

\begin{itemize}
\item \emph{Spectral determination of masses}
\[
\text{Masses} = \sqrt{\operatorname{spec}(H_f)}.
\]

\item \emph{Orientational determination of mixing}
\[
\text{Mixing} = U_f^\dagger U_{f'}.
\]
\end{itemize}

The spectral triple determines all mass eigenvalues, while the alignment mechanism determines the otherwise undetermined orientation between degenerate fibers.

This two-step geometric picture unifies the understanding of flavor: masses arise from spectra, mixing from orientation.

\subsection{Implications for the Geometry of the Standard Model}

These results point toward a deeper geometric interpretation of the Standard Model:

\begin{itemize}
\item Degenerate fibers in multiplicity spaces encode internal symmetries not fixed by the spectral triple.
\item Relative orientations of such fibers carry physical meaning.
\item Minimal relational operators can lift degeneracies and yield definite physical predictions.
\item The full flavor structure emerges not from additional algebraic input, but from the geometry of sector interrelations.
\end{itemize}

This perspective suggests possible extensions:

\begin{itemize}
\item Connections with renormalization flows or dynamical alignment processes;
\item Generalizations to other multiplicity structures in extended spectral triples;
\item Possible geometric origins of CP phases.
\end{itemize}

These directions provide fertile ground for future research.

\subsection{Summary}

The alignment mechanism presented here transforms the remaining ambiguity in the PMNS matrix into a mathematically well-posed variational problem, whose solution yields the physical atmospheric mixing angle.

This mechanism:
\begin{itemize}
\item respects all NCG axioms;
\item introduces no new degrees of freedom;
\item is uniquely determined by geometric constraints;
\item preserves all existing flavor data;
\item and resolves the final undetermined parameter of the lepton mixing matrix.
\end{itemize}

It thereby completes the geometric interpretation of flavor within the noncommutative Standard Model, showing that the full structure of lepton mixing emerges from the interplay between spectra and orientation in generation space.


\section{Conclusion}

In this work we have identified and resolved a structural limitation in the finite spectral triple formulation of the Standard Model related to fermion mixing. While the noncommutative geometric framework determines fermion mass spectra through the eigenvalues of the sector Hermitian operators
\[
H_f = Y_f Y_f^\dagger,
\]
it does not, in the presence of degeneracies, determine the relative orientation of the corresponding eigenspaces. When the effective neutrino Hermitian operator admits a two-dimensional eigenspace-whether exactly or approximately-this results in a residual internal $\mathrm{SO}(2)$ freedom that is invisible to the spectral action and compatible with all axioms of the spectral triple, yet acts nontrivially on the PMNS matrix by shifting the atmospheric mixing angle $\theta_{23}$.

From a broader perspective, the results presented here clarify the scope of spectral geometry. While the spectral action determines eigenvalue data and constrains mass hierarchies, it does not, in the presence of degeneracies, determine relative orientations within internal multiplicity spaces. This limitation is not a defect of the framework, but a structural feature. The alignment mechanism proposed here makes this feature explicit and provides a minimal relational completion, without modifying the axioms or extending the spectral triple.

Our first main result establishes that the residual $\mu$-$\tau$ rotational degree of freedom is the geometric origin of the undetermined atmospheric mixing angle in spectral triple formulations. This explains why $\theta_{23}$ remains unconstrained within the standard noncommutative geometric framework: it corresponds to an internal orientation in generation space that is not fixed by spectral data alone.

Our second main result identifies the alignment operator
\[
A(\alpha,\beta)=\alpha H_e + \beta H_\nu,
\]
as the unique minimal relational operator-compatible with the order-zero condition, first-order condition, real structure, grading, and spectral stability-that is capable of resolving this residual $\mathrm{SO}(2)$ freedom by selecting a preferred basis in the $\mu$-$\tau$ plane. The operator depends solely on geometric data already present in the spectral triple and introduces no new algebraic structure or Dirac components.

To quantify the relative misalignment between sector operators, we introduced the alignment energy functional
\[
E(\theta)
= \big\| [H_e, R(\theta)^\dagger H_\nu R(\theta)] \big\|^2,
\]
which measures the failure of the charged-lepton and neutrino Hermitian operators to be simultaneously diagonal at angle $\theta$. This functional is analytic, periodic, and depends only on the relative geometry of $H_e$ and $H_\nu$. We proved that it possesses a unique global minimum at an angle $\theta$, corresponding to the preferred internal orientation selected by the alignment mechanism.

At this angle, the PMNS matrix takes the aligned form
\[
U_{\mathrm{PMNS}}^{(\mathrm{aligned})}
= U_e^\dagger U_\nu R(\theta),
\]
and the atmospheric mixing angle $\theta_{23}$ is uniquely fixed.
Crucially, all remaining flavor data-charged-lepton masses, quark masses, neutrino mass-squared differences, CKM mixing parameters, and the solar and reactor mixing angles-remain entirely unchanged. The alignment mechanism is therefore selective and surgical: it resolves only the single geometric ambiguity identified in Section~3 and does not alter any established phenomenological successes of the Standard Model.

We emphasize that claims of uniqueness in this work are meant in a restricted sense: among linear, relational operators constructed from existing sector Hermitians and compatible with the axioms, the alignment operator represents the simplest choice. More elaborate constructions are logically possible but would introduce additional structure beyond the minimal scope considered here.

Conceptually, this work provides a unified geometric interpretation of flavor within the noncommutative Standard Model. Mass spectra arise from the eigenvalues of sector Hermitian operators, while mixing angles arise from the relative orientation of their eigenbases in generation space. Degeneracies correspond to internal symmetries, and the alignment operator supplies the minimal relational structure required to extract physical predictions when such symmetries are present.

Several natural directions follow from this analysis.
The alignment principle may extend to CP-violating phases or to other sectors in which internal multiplicities arise, suggesting the possibility of a fully geometric interpretation of CP violation. The relationship between alignment, renormalization-group stability, and dynamical flows also merits further study. More generally, the role of multiplicity spaces in extended or alternative almost-commutative geometries presents a promising avenue for future work.

In summary, the alignment mechanism presented here completes the geometric determination of flavor up to internal orientation in the noncommutative Standard Model. By resolving the residual $\mu$-$\tau$ freedom left undetermined by the spectral triple, it yields a uniquely determined atmospheric mixing angle and demonstrates how the full lepton mixing structure can be extracted from the intrinsic geometry of fermion multiplets. This establishes a new level of conceptual coherence within the noncommutative geometric framework and opens the door to further geometric insights into the structure of particle physics.

\appendix

\section{Appendix}

GitHub Repository: \url{https://github.com/ChazzCoin/Yukawa-Emergence.git}


\end{document}
